\section{Graphical languages for quantum circuits}\label{sec:circuits}

The formal language is categorical, but its role is practical: it fixes the syntax for the usual
boxes-and-wires calculus.  In this setting, generic wiring is treated structurally, so the
displayed rule sets can focus on fragment-specific generators and equations.

\subsection{Graphical languages}

We work with the one-object-per-wire formalisms PROs and PROPs in the sense of
Joyal--Street~\cite{joyal_geometry_1991, selinger_survey_2010}.
% The intended operational reading is standard: objects are numbers of wires, sequential composition
% is plugging, and tensor is side-by-side juxtaposition.  
We adopt strictness throughout: associativity
and unit constraints are treated as literal equalities, so parentheses can be omitted.

\begin{definition}
A \emph{PRO} is a strict monoidal category \(\cat{P}\) whose objects are the
natural numbers \(n\in\N\) and whose tensor on objects is addition
\(n\otimes m \defeq n+m\) with unit \(0\).  We write \(\cat{P}(n,m)\) for morphisms
\(n\to m\), with sequential composition \(\circ\) and parallel composition \(\otimes\),
subject to the usual strict monoidal axioms
(\cref{eq:seqidentity,eq:seqassociativity,eq:parassociativity,eq:paridentity,eq:interchange}).
\end{definition}
% Intuitively, a PRO is a typed circuit calculus with \(n\) input wires and \(m\) output wires, 
% where
% \(\circ\) is plugging and \(\otimes\) is juxtaposition, but where wire permutations are not available as
% structural morphisms.

We write \(f:n\to m\) when \(f\in\cat{P}(n,m)\); \(n\) and \(m\) are the numbers of input
and output wires of~\(f\).  In particular \(\id_n=\id_1^{\otimes n}\) for all \(n\in\N\).

\begin{definition}
A \emph{PROP} is a PRO \(\cat{P}\) equipped with a symmetric monoidal structure,
i.e. symmetry isomorphisms \(\sigma_{n,m}:n+m\to m+n\) \((n,m\in\N)\) generated by
the basic swap \(\sigma\defeq\sigma_{1,1}\) and satisfying the axioms
\eqref{eq:involution}-\eqref{eq:naturality} in \cref{fig:coherenceprop}.  We
write \(\sigma_n\defeq\sigma_{n,1}:n+1\to 1+n\) for swapping one wire with a block
of \(n\) wires.
\end{definition}
In a PROP, swaps and permutations are part of the ambient structure, and the resulting
string-diagram representation draws a morphism \(f:n\to m\) as a box with \(n\) incoming wires and
\(m\) outgoing wires.
% (and hence will not be
% counted as fragment-specific axioms later on).
Sequential composition is drawn by stacking boxes, tensor by placing them side by side, and
\(\sigma_{n,m}\) by crossing a block of \(n\) wires with a block of \(m\) wires.  With these conventions,
the coherence axioms are exactly the diagrammatic moves that make deformation sound.
For example, if \(g:1\to1\) and \(h:2\to2\) are gates, then \(g\otimes h:3\to3\) is their
side-by-side circuit, and composing with \(\sigma_{1,2}\) just moves the first output wire past the
other two.  Such crossings are part of the syntax before any fragment-specific rule is added.

\begin{center}
  \begin{tabular}{c@{\hskip 1.4em}c@{\hskip 1.4em}c@{\hskip 1.4em}c@{\hskip 1.4em}c@{\hskip 1.4em}c@{\hskip 1.4em}c@{\hskip 1.4em}c@{\hskip 1.4em}}
    \(f\) & \(\id_0\) & \(\id_1\) & \(\id_n\) &
    \(\sigma\) & \(\sigma_{n}\) & \(g\circ f\) & \(f\otimes g\) \\
    \tikzfig{category/f} & \tikzfig{category/empty} &
    \tikzfig{category/id} & \tikzfig{category/idn} &
    \tikzfig{gates/SWAP} & \tikzfig{category/perm} &
    \tikzfig{category/fg} & \tikzfig{category/fparg}
  \end{tabular}
\end{center}

% The coherence axioms for PROs and PROPs, used throughout the paper, are
% summarised in \cref{fig:coherenceprop}.  
Equalities obtained solely by these coherence moves, including permutations, are
\emph{structural} and may be omitted from diagrammatic arguments.

\captionsetup[subfigure]{justification=centering}
%\newpage
\begin{figure}[ht]
  \centering
  \begin{subfigure}[t]{0.43\textwidth}
    \centering
    \scalebox{0.9}{\fbox{\tikzfig{category/seqidentity}}}
    \caption{\(f\circ\id_n = f = \id_n\circ f\)}
    \label{eq:seqidentity}
  \end{subfigure}%  
  \begin{subfigure}[t]{0.43\textwidth}
    \centering
    \scalebox{0.9}{\fbox{\tikzfig{category/seqassociativity}}}
    \caption{\(h\circ(g\circ f) = (h\circ g)\circ f\)}
    \label{eq:seqassociativity}
  \end{subfigure}
  %\vspace{1ex}
  \begin{subfigure}[t]{0.30\textwidth}
    \centering
    \scalebox{0.9}{\fbox{\tikzfig{category/parassociativity}}}
    \caption{\((f\!\otimes\! g)\!\otimes\! h \!=\! f\!\otimes\!(g\!\otimes\! h)\)}
    \label{eq:parassociativity}
  \end{subfigure}%   
  \begin{subfigure}[t]{0.45\textwidth}
    \centering
    \scalebox{0.9}{\fbox{\tikzfig{category/interchange}}}
    \caption{\((g_1\!\circ\! f_1)\!\otimes\!(g_2\!\circ\! f_2)
         \!=\! (g_1\!\otimes\! g_2)\!\circ\!(f_1\!\otimes\! f_2)\)}
    \label{eq:interchange}
  \end{subfigure}
  
  \vspace{1ex}

  \begin{subfigure}[t]{0.38\textwidth}
    \centering
    \scalebox{0.9}{\fbox{\tikzfig{category/paridentity}}}
    \caption{\(f \otimes \id_0 = f = \id_0 \otimes f\)}
    \label{eq:paridentity}
  \end{subfigure}% 
  \begin{subfigure}[t]{0.26\textwidth}
    \centering
    \scalebox{0.9}{\fbox{\tikzfig{category/involution}}}
    \caption{\(\sigma\circ\sigma = \id_2\)}
    \label{eq:involution}
  \end{subfigure}%  
  % empty centre cell
  % \begin{subfigure}[t]{0.10\textwidth}
  %   \mbox{}
  % \end{subfigure}%  
  % \hfill
  \begin{subfigure}[t]{0.35\textwidth}
    \centering
    \scalebox{0.9}{\fbox{\tikzfig{category/naturality}}}
    \caption{\(\sigma_m\circ(f\otimes\id_1) = (\id_1\otimes f)\circ\sigma_n\) for \(f:n\to m\)}
    \label{eq:naturality}
  \end{subfigure}
  \caption{Coherence laws for PROs and PROPs.}
  \label{fig:coherenceprop}
\end{figure}

\subsection{Free PROPs and presentations}\label{free-prop}

Circuit fragments are specified by \emph{generators and relations}.  Generators represent primitive
gates with fixed arity, and relations are the rewrite rules used for equational reasoning.

\begin{definition}\label{def:signature}
A \emph{(PROP) signature} is a set \(\Sigma\) equipped with arity functions
\(\dom,\cod:\Sigma\to\N\).
An element \(g\in\Sigma\) is written as a generating morphism \(g:\dom(g)\to\cod(g)\).
\end{definition}

\begin{definition}\label{def:free-prop}
Let \(\Sigma\) be a signature.  The \emph{free PROP} on \(\Sigma\), written \(\cat{P}_\Sigma\), is the
PROP whose morphisms are string diagrams generated from:
\begin{itemize}
\item the non-structural generators \(g\in\Sigma\) (with their prescribed arities), and
\item the structural morphisms of a PROP (identities and symmetries),
\end{itemize}
quotiented by PROP coherence (i.e.\ the axioms of \cref{fig:coherenceprop}).
\end{definition}

\begin{definition}\label{def:presented-prop}
Let \(\Sigma\) be a signature and let \(\mathcal{R}\) be a set of well-typed equations
\(L=R\) between morphisms of the free PROP \(\cat{P}_\Sigma\).
The \emph{presented PROP} \(\cat{P}_\Sigma/\mathcal{R}\) is the quotient of \(\cat{P}_\Sigma\)
by the smallest congruence that contains \(\mathcal{R}\) and is closed under \(\circ\), \(\otimes\),
and symmetry.
\end{definition}

This quotient is the \emph{graphical language} used in the rest of the paper: diagrams are equal exactly when their
equality is derivable from PROP coherence together with the chosen relations.

% Equivalences in \(\cat{P}_\Sigma/\mathcal{R}\) are derived from the equations in
% \(\mathcal{R}\) and the coherence axioms of \cref{fig:coherenceprop} using the
% usual equational rules (reflexivity, symmetry, transitivity and congruence for
% \(\circ\) and \(\otimes\)); we do not spell them out explicitly.

% \begin{remark}\label{rem:structural-vs-derivable}
% We distinguish:
% \begin{itemize}
% \item \emph{structural equality} in \(\cat{P}_\Sigma\), written \(C_1\equiv_\Sigma C_2\),
%       meaning equality using only PROP coherence (associativity/unit laws and symmetries); and
% \item \emph{derivability} in a presented PROP \(\cat{P}_\Sigma/\mathcal{R}\), written
%       \(\cat{P}_\Sigma/\mathcal{R}\vdash C_1=C_2\), meaning equality derivable from
%       \(\mathcal{R}\) using equational reasoning (plus structural moves).
% \end{itemize}
% \end{remark}

\subsection{Monoidal functors and completeness}

An interpretation sends the presented syntax to its intended operations; completeness is the
assertion that every semantic equality in that image can be recovered as a derivation in the
quotient.

\begin{definition}\label{def:prop-morphism}
A \emph{PROP morphism} is a strict symmetric monoidal functor
\(F:\cat{P}_1\to\cat{P}_2\) between PROPs that is the identity on objects.
Equivalently, \(F\) preserves identities, symmetries, sequential composition and
tensor.
\end{definition}

\begin{definition}\label{def:subprop}
Let \(\cat{P}\) be a PROP.
A \emph{subPROP} of \(\cat{P}\) is a wide strict symmetric monoidal subcategory
\(\cat{Q}\subseteq\cat{P}\).
Since the objects of a PROP are fixed to be \(\mathbb{N}\), this amounts to specifying
subsets
\(
  \cat{Q}(n,m)\subseteq \cat{P}(n,m)\;(n,m\in\mathbb{N})
\)
such that \(\id_n\in \cat{Q}(n,n)\) for all \(n\); \(\sigma_{n,m}\in \cat{Q}(n+m,m+n)\) for all \(n,m\); and
\(\cat{Q}\) is closed under \(\circ\) and \(\otimes\).

In this situation, the inclusion functor \(\cat{Q}\hookrightarrow \cat{P}\) is a faithful PROP morphism.
\end{definition}

\begin{definition}\label{def:interpretation-functor}
Let \(\cat{P}_\Sigma/\mathcal{R}\) be a presented PROP and let \(\cat{C}\) be a
symmetric monoidal category. An \emph{interpretation functor} is a strict
symmetric monoidal functor
\(
  \interp{\cdot}:\cat{P}_\Sigma\to\cat{C}
\)
specified on generators \(g\in\Sigma\) and extended monoidally to all diagrams.

If every equation in \(\mathcal{R}\) holds in \(\cat{C}\), then
\(\interp{\cdot}\) factors uniquely through the quotient, yielding a functor
\(
  \interp{\cdot}:\cat{P}_\Sigma/\mathcal{R}\to\cat{C}.
\)
\end{definition}

\begin{definition}\label{def:semantic-completeness}
Let \(\interp{\cdot}:\cat{P}_\Sigma/\mathcal{R}\to\cat{C}\) be such an
interpretation.  We say that the graphical language \(\cat{P}_\Sigma/\mathcal{R}\)
is \emph{complete} for this semantics if, for all circuits \(C_1,C_2\), \(\interp{C_1}=\interp{C_2}\) implies
\(\cat{P}_\Sigma/\mathcal{R}\vdash C_1=C_2\).  Equivalently, the interpretation
functor is faithful.

We will also use a bounded variant. Given \(k\in\mathbb{N}\), we say that
\(\cat{P}_\Sigma/\mathcal{R}\) is \emph{complete up to \(k\) wires} if the
implication holds whenever \(C_1,C_2:n\to m\) with \(n,m\le k\).  Equivalently, the
interpretation functor is faithful on each hom-set
\(\cat{P}_\Sigma/\mathcal{R}(n,m)\) for \(n,m\le k\).
\end{definition}
For equational reasoning, completeness says that semantic equality can always be witnessed by an explicit
diagrammatic derivation from the chosen relations; the bounded form restricts this requirement to
circuits of limited arity.

When \(\cat{C}\) is, by definition, the symmetric monoidal subcategory of an
ambient category generated by the images of the generators, the interpretation
functor is also full; in that case \(\cat{P}_\Sigma/\mathcal{R}\) is actually
isomorphic to the semantic subPROP~\(\cat{C}\).

Completeness theorems stated for PROs can still be used in this PROP setting through a single
lifting principle.

\begin{lemma}[PRO-to-PROP lifting]\label{lem:fromprotoprop}\label{rem:fromprotoprop}
Let \(\cat{P}/\mathcal{R}\) be a complete presented PRO with an interpretation into a \emph{symmetric} monoidal category \(\cat{C}\). 
Assume there exists a PRO circuit
\(\tau:2\to 2\) such that \(\interp{\tau}\) is the symmetry \(\sigma_{1,1}\) of
\(\cat{C}\).
Form the free PROP on the same generators, and quotient it by \(\mathcal{R}\) together with the single equation
\(\sigma_{1,1}=\tau\).
Then the resulting PROP presentation is complete for the same semantics.
\end{lemma}
\begin{proof}
The equation \(\sigma_{1,1}=\tau\) rewrites each structural swap between adjacent wires as a source PRO circuit.
By the symmetric monoidal axioms of a PROP, every structural permutation is generated by such adjacent swaps,
so each PROP diagram \(D\) has a translated PRO diagram \(T(D)\) obtained by replacing structural swaps by copies of
\(\tau\) and leaving the ordinary generators unchanged.  The interpretation satisfies
\(\interp{D}=\interp{T(D)}\) because \(\interp{\tau}=\sigma_{1,1}\).  If two PROP diagrams have the same semantics,
their translations have the same semantics as PRO diagrams, hence are equal by PRO completeness.  Replacing
each occurrence of \(\tau\) by \(\sigma_{1,1}\) then gives the original PROP equality.
\end{proof}
\section{Quantum circuits and their relations}
\label{sec:circuits-relations}

The transfer and separation arguments use one Hilbert-space semantic setup: five qubit fragments
share a signature, while qutrit Clifford uses its own presented subPROP.

Let \(\cat{FdHilb}\) be the symmetric monoidal category of finite-dimensional
complex Hilbert spaces and linear maps, with tensor product~\(\otimes\), tensor
unit~\(\C\), and symmetry given by swapping tensor factors.  By Mac~Lane's
coherence theorem \cite{Saunders} we may regard \(\cat{FdHilb}\) as strict monoidal.

\begin{definition}\label{def:unitary}
Let \(H,K\) be finite-dimensional Hilbert spaces. A linear map \(U:H\to K\) is \emph{unitary} if
\(
  U^\dagger U = \id_H
\) and \(
  UU^\dagger = \id_K,
\)
where \(U^\dagger\) denotes the Hilbert adjoint of \(U\).
We write \(\mathrm{U}(H,K)\) for the set of unitary maps \(H\to K\), and \(\mathrm{U}(H)\coloneqq \mathrm{U}(H,H)\).
\end{definition}

\begin{definition}\label{def:qubit-qutrit-props}
For \(d\ge 2\), let \(\cat{Qudit}_d\) be the PROP with objects \(\N\) and
hom-sets \(\cat{Qudit}_d(n,m)\coloneqq
\mathrm{U}\big((\C^d)^{\otimes n},(\C^d)^{\otimes m}\big)\), with composition,
identities, tensor, and symmetry inherited from \(\cat{FdHilb}\) viewed as
strict monoidal.  The qubit and qutrit PROPs are
\(\cat{Qubit}\coloneqq\cat{Qudit}_2\) and
\(\cat{Qutrit}\coloneqq\cat{Qudit}_3\).
\end{definition}
For \(d\ge2\), \(\cat{Qudit}_d\) is endomorphism-only: there are no unitaries
\(n\to m\) for \(n\neq m\).
For a set \(G\) of morphisms in a PROP, \(\langle G\rangle\) denotes the smallest subPROP
containing \(G\) and all structural morphisms.

\subsection{Qubit circuit fragments as presented subPROPs}

The five qubit fragments share a common qubit signature and unitary interpretation.

\begin{definition}
The \emph{common qubit signature} \(\Sigma^{(2)}\) has non-structural generators
\(\galpha,\gminus:0\to0\),
\(\gH,\gZ,\gS,\gT,\gX:1\to1\), and \(\gCNOT,\gCS:2\to2\).
We write \(\cat{P}^{(2)}\defeq\cat{P}_{\Sigma^{(2)}}\) for the free PROP on
\(\Sigma^{(2)}\).  Structural morphisms \(\id_n\) and \(\sigma_{n,m}\) come from the
ambient PROP structure.
\end{definition}

\begin{remark}\label{rem:arity-preservation}
All signatures used in this section are endomorphism-only: each generator has type \(n\to n\).
Hence every well-typed fragment circuit, and every equation between such circuits, preserves the
wire count; this is the arity fact used in \cref{sec:minimality}.
\end{remark}

We use the standard Dirac notation for computational basis states: for \(x\in\{0,1\}\) we write \(\ket{x}\in\C^2\),
and for \(x,y\in\{0,1\}\) we abbreviate \(\ket{x,y}\coloneqq \ket{x}\otimes\ket{y}\) (similarly,
\(\ket{x_1,\ldots,x_n}\coloneqq \ket{x_1}\otimes\cdots\otimes\ket{x_n}\)).

By the universal property of the free PROP, there is a unique strict
symmetric monoidal functor
\(
  \interp{\cdot}^{(2)}:\cat{P}^{(2)}\to\cat{Qubit}
\)
sending:
\begin{equation*}
\begin{array}{l l}
 \interp{\galpha}^{(2)} : 1\longmapsto e^{i\pi/4} &
 \interp{\gH}^{(2)} : \ket{x}\longmapsto
     \tfrac1{\sqrt2}\!\sum_{k=0}^1(-1)^{xk}\ket{k}\\
 \interp{\gminus}^{(2)} : 1\longmapsto -1 &
 \interp{\gS}^{(2)} : \ket{x}\longmapsto i^x\ket{x} \\
 \interp{\gX}^{(2)} : \ket{x}\longmapsto\ket{1-x} &
 \interp{\gZ}^{(2)} : \ket{x}\longmapsto(-1)^x\ket{x} \\
 \interp{\gCNOT}^{(2)} : \ket{x,y}\longmapsto\ket{x,x\oplus y} &
 \interp{\gT}^{(2)} : \ket{x}\longmapsto e^{i\pi x/4}\ket{x} \\
 \interp{\gCS}^{(2)} : \ket{x,y}\longmapsto i^{xy}\ket{x,y}
\end{array}
\end{equation*}
for \(x,y\in\{0,1\}\) and addition \(\oplus\) modulo~2.

Each qubit fragment is specified by a sub-signature of \(\Sigma^{(2)}\) and equations on the free
PROP it generates; the common unitary interpretation then restricts to this syntax.  The index
\(\bullet\) records the chosen fragment, while \(\Sigma_\bullet\), \(\cat{P}_\bullet\),
\(\QC_\bullet\), \(\QCirc_\bullet\), \(\cat{Qubit}_\bullet\), and
\(\interp{\cdot}_\bullet\) denote its gates, free PROP, equations, presented PROP, generated
semantic subPROP, and quotient interpretation, respectively.

\begin{definition}\label{def:qubit-fragment-scheme}
Let
\(
  \bullet \in \{\mathit{Cliff}, \mathit{RCliff}, \mathit{CliffT}, \mathit{CliffCS}, \mathit{CNOTdihe}\}
\)
range over the five qubit fragments considered in this paper.  Given a sub-signature \(\Sigma_\bullet\subseteq\Sigma^{(2)}\), let
\(\cat{P}_\bullet\defeq\cat{P}_{\Sigma_\bullet}\) be the free PROP on \(\Sigma_\bullet\).  If
\(\QC_\bullet\) is a finite set of well-typed equations in \(\cat{P}_\bullet\), define the presented
PROP
\(
  \QCirc_\bullet\defeq\cat{P}_\bullet/\QC_\bullet.
\)
The semantic target is the subPROP
\(
  \cat{Qubit}_\bullet
  \defeq
  \big\langle\,\interp{g}^{(2)} \mid g\in\Sigma_\bullet\,\big\rangle
  \subseteq
  \cat{Qubit},
\)
generated by the interpretations of the gates in \(\Sigma_\bullet\).  When \(\QC_\bullet\) is sound
for the restricted common interpretation, the quotient carries the induced strict symmetric
monoidal functor
\(
  \interp{\cdot}_\bullet:\QCirc_\bullet\to\cat{Qubit}_\bullet.
\)
We call such a presentation complete when this functor is faithful.
\end{definition}

\begin{definition}\label{def:qubit-fragments}
The five qubit fragments treated in this paper are specified by these sub-signatures and rule sets:
\begin{center}
\begin{tabular}{@{}l p{0.38\linewidth} p{0.38\linewidth}@{}}
\textbf{Fragment} & \(\Sigma_\bullet\subseteq\Sigma^{(2)}\) & \textbf{Rule set \(\QC_\bullet\)} \\
\hline
Clifford
& \(\Sigma_{\mathit{Cliff}}\defeq\{\galpha,\gH,\gS,\gCNOT\}\)
& \(\QC_{\mathit{Cliff}}\defeq\QCclifford\) (\cref{fig:qubitcliffordaxioms})
\\
Real Clifford
& \(\Sigma_{\mathit{RCliff}}\defeq\{\gminus,\gH,\gZ,\gCNOT\}\)
& \(\QC_{\mathit{RCliff}}\defeq\QCrealclifford\) (\cref{fig:qubitrealcliffordaxioms})
\\
Clifford\(+T\)
& \(\Sigma_{\mathit{CliffT}}\defeq\{\galpha,\gH,\gT,\gCNOT\}\)
& \(\QC_{\mathit{CliffT}}\defeq\QCcliffordplust\) (\cref{fig:qubitcliffordplustaxioms})
\\
Clifford\(+CS\)
& \(\Sigma_{\mathit{CliffCS}}\defeq\{\galpha,\gH,\gS,\gCS\}\)
& \(\QC_{\mathit{CliffCS}}\defeq\QCcliffordpluscs\) (\cref{fig:qubitcliffordpluscsaxioms})
\\
CNOT-dihedral
& \(\Sigma_{\mathit{CNOTdihe}}\defeq\{\galpha,\gX,\gT,\gCNOT\}\)
& \(\QC_{\mathit{CNOTdihe}}\defeq\QCdihedral\) (\cref{fig:qubitdihedralaxioms})
\end{tabular}
\end{center}
\end{definition}

\subsection{The qutrit Clifford fragment}
\label{subsec:qutrit-clifford-fragment}

For qutrits, the paper uses only the Clifford fragment, denoted \(\mathit{Cliff3}\), whose
signature is separate from the common qubit signature; we use the computational basis
\(\{\ket{0},\ket{1},\ket{2}\}\) of \(\C^3\) and the induced tensor-product basis.

\begin{definition}\label{def:qutrit-clifford-fragment}
The \emph{qutrit Clifford} fragment is generated by the signature
\(\Sigma_{\mathit{Cliff3}}\defeq
\{\gww:0\to0,\gH,\gS:1\to1,\gCNOT:2\to2\}\).

Let \(\cat{P}_{\mathit{Cliff3}}\defeq\cat{P}_{\Sigma_{\mathit{Cliff3}}}\) be the free PROP on
\(\Sigma_{\mathit{Cliff3}}\) and let \(\QC_{\mathit{Cliff3}}\defeq\QCqutritclifford\) be the rule
set given in \cref{fig:qutritcliffordaxioms}.  The presented PROP is
\(
  \QCirc_{\mathit{Cliff3}}\defeq\cat{P}_{\mathit{Cliff3}}/\QC_{\mathit{Cliff3}}.
\)
\end{definition}

By the universal property of the free PROP, the assignments
\begin{equation*}
\begin{array}{l l}
 \interp{\gww}_{\mathit{Cliff3}} : 1\longmapsto e^{i\pi/6} & 
 \interp{\gH}_{\mathit{Cliff3}} : \ket{x}\longmapsto 
   \tfrac1{\sqrt3}\!\sum_{k=0}^2 e^{2\pi i\,xk/3}\ket{k} \\
 \interp{\gCNOT}_{\mathit{Cliff3}} : \ket{x,y}\longmapsto \ket{x,(x+y)\bmod 3} &
 \interp{\gS}_{\mathit{Cliff3}} : \ket{x}\longmapsto
   e^{i\pi x(x-1)/3}\ket{x} 
\end{array}
\end{equation*}
for \(x,y\in\{0,1,2\}\) determine a strict symmetric monoidal functor
\(\cat{P}_{\mathit{Cliff3}}\to\cat{Qutrit}\).  Let
\(\cat{Qutrit}_{\mathit{Cliff3}}\defeq\langle\,\interp{g}_{\mathit{Cliff3}}\mid g\in\Sigma_{\mathit{Cliff3}}\,\rangle\)
be the semantic subPROP generated by the qutrit Clifford gates.  Since the equations of
\(\QC_{\mathit{Cliff3}}\) are sound for this interpretation, the free-PROP interpretation factors
through the quotient and yields
\(
  \interp{\cdot}_{\mathit{Cliff3}}:\QCirc_{\mathit{Cliff3}}\longrightarrow\cat{Qutrit}_{\mathit{Cliff3}}.
\)

\subsection{Rule-set presentations}

Each displayed equation is an arity-preserving schema; gates outside the corresponding fragment
signature are shortcuts expanded as in \cref{fig:quantum-circuit-shortcuts}.

%%%%%%%%%%%%%%%%%%%%%%%%%%%%%%%%%%%%%%%%%%%%%%%%%%%%%%%%%%%%%%%%%%%%%%%%%%%%%%%%%%%%%
%%%%%%%%%%%%%%%%%%%%%%%%%%%%%%%%%%% QUBIT CLIFFORD %%%%%%%%%%%%%%%%%%%%%%%%%%%%%%%%%%
%%%%%%%%%%%%%%%%%%%%%%%%%%%%%%%%%%%%%%%%%%%%%%%%%%%%%%%%%%%%%%%%%%%%%%%%%%%%%%%%%%%%%
\begin{figure}[ht]
    \centering
    \fbox{\begin{minipage}{0.975\linewidth}
    \centering
    \vspace{-1em}
    \begin{minipage}[c]{0.22\linewidth}\centering
      \begin{equation}\tag{\omegarulabel{8}}\label{w8}
        \scalebox{0.9}{\tikzfig{gates/wc}}^{\otimes 8} = \scalebox{0.9}{\tikzfig{gates/empty}}
      \end{equation}
    \end{minipage}%
    \begin{minipage}[c]{0.24\linewidth}\centering
      \begin{equation}\tag{\rulabel{H^2}}\label{H2}
        \scalebox{0.9}{\tikzfig{cliffordaxioms/HH}} = \scalebox{0.9}{\tikzfig{gates/Id}}
      \end{equation}
    \end{minipage}%
    \begin{minipage}[c]{0.32\linewidth}\centering
      \begin{equation}\tag{\rulabel{S^4}}\label{S4}
        \scalebox{0.9}{\tikzfig{cliffordaxioms/SSSS}} = \scalebox{0.9}{\tikzfig{gates/Id}}
      \end{equation}
    \end{minipage}\vspace{-0.6em}
    \begin{minipage}[c]{0.54\linewidth}\centering
      \begin{equation}\tag{E}\label{E}
        \scalebox{0.9}{\tikzfig{cliffordaxioms/HSH}} = \scalebox{0.9}{\tikzfig{cliffordaxioms/wSdaggerHSdagger}}
      \end{equation}
    \end{minipage}%
    \begin{minipage}[c]{0.32\linewidth}\centering
      \begin{equation}\tag{CPh}\label{Cs}
        \scalebox{0.9}{\tikzfig{cliffordaxioms/CNOTSCNOT}} = \scalebox{0.9}{\tikzfig{cliffordaxioms/S}}
      \end{equation}
    \end{minipage}\vspace{-0.6em}
    \begin{minipage}[c]{0.29\linewidth}\centering
      \begin{equation}\tag{B}\label{B}
        \scalebox{0.9}{\tikzfig{cliffordaxioms/CNOTNOTC}} = \scalebox{0.9}{\tikzfig{cliffordaxioms/SWAPCNOT}}
      \end{equation}
    \end{minipage}%
    \begin{minipage}[c]{0.47\linewidth}\centering
      \begin{equation}\tag{CZ}\label{CZ}
        \scalebox{0.9}{\tikzfig{cliffordaxioms/HCNOTH}} = \scalebox{0.9}{\tikzfig{cliffordaxioms/CZ}}
      \end{equation}
    \end{minipage}\vspace{-0.6em}
    \begin{minipage}[c]{0.29\linewidth}\centering
      \begin{equation}\tag{I}\label{Inew}
        \scalebox{0.9}{\tikzfig{qutritcliffordaxioms/CNOT23CNOT12CNOT13}} = \scalebox{0.9}{\tikzfig{qutritcliffordaxioms/CNOT12CNOT23}}
      \end{equation}
    \end{minipage}
    \end{minipage}}
    \caption{Qubit Clifford presentation \(\QCclifford\), with the scalar, one-qubit, controlled-phase, swap-decomposition, controlled-\(Z\), and three-wire interaction laws used in the completeness and minimality arguments.}
    \label{fig:qubitcliffordaxioms}
\end{figure}
%%%%%%%%%%%%%%%%%%%%%%%%%%%%%%%%%%%%%%%%%%%%%%%%%%%%%%%%%%%%%%%%%%%%%%%%%%%%%%%%%%%%%%%%%%
%%%%%%%%%%%%%%%%%%%%%%%%%%%%%%%%%%% QUBIT REAL CLIFFORD %%%%%%%%%%%%%%%%%%%%%%%%%%%%%%%%%%
%%%%%%%%%%%%%%%%%%%%%%%%%%%%%%%%%%%%%%%%%%%%%%%%%%%%%%%%%%%%%%%%%%%%%%%%%%%%%%%%%%%%%%%%%%
\begin{figure}[ht]
    \centering
    \fbox{\begin{minipage}{0.975\linewidth}
        \centering
        \vspace{-1em}
        %%%%%% 0-wire & 1-wire %%%%%%
        \begin{minipage}[c]{0.22\linewidth}\centering
            \begin{equation}\tag{\rulabel{-^2}}\label{realminus2}
                \scalebox{0.9}{\tikzfig{gates/minus}}^{\otimes 2} = \scalebox{0.9}{\tikzfig{gates/empty}}
            \end{equation}
        \end{minipage}
        \begin{minipage}[c]{0.24\linewidth}\centering
            \begin{equation}\tag{\rulabel{H^2}}\label{realH2}
                \scalebox{0.9}{\tikzfig{realcliffordaxioms/H2}} = \scalebox{0.9}{\tikzfig{gates/Id}}
            \end{equation}
        \end{minipage}
        \begin{minipage}[c]{0.24\linewidth}\centering
            \begin{equation}\tag{\rulabel{Z^2}}\label{realZ2}
                \scalebox{0.9}{\tikzfig{realcliffordaxioms/Z2}} = \scalebox{0.9}{\tikzfig{gates/Id}}
            \end{equation}
        \end{minipage}\vspace{-0.6em}
        \begin{minipage}[c]{0.50\linewidth}\centering
            \begin{equation}\tag{F}\label{realF}
                \scalebox{0.9}{\tikzfig{realcliffordaxioms/ZHZH}} = \scalebox{0.9}{\tikzfig{realcliffordaxioms/minusHZHZ}}
            \end{equation}
        \end{minipage}
        %%%%%% 2-wire %%%%%%
        \begin{minipage}[c]{0.30\linewidth}\centering
            \begin{equation}\tag{\rulabel{CX^2}}\label{realCX2}
                \scalebox{0.9}{\tikzfig{realcliffordaxioms/CNOTCNOT}} = \scalebox{0.9}{\tikzfig{realcliffordaxioms/II}}
            \end{equation}
        \end{minipage}\vspace{-0.6em}
        \begin{minipage}[c]{0.29\linewidth}\centering
            \begin{equation}\tag{B}\label{realB}
                \scalebox{0.9}{\tikzfig{cliffordaxioms/CNOTNOTC}} = \scalebox{0.9}{\tikzfig{cliffordaxioms/SWAPCNOT}}
            \end{equation}
        \end{minipage}
        \begin{minipage}[c]{0.35\linewidth}\centering
            \begin{equation}\tag{ZC}\label{realwC}
                \scalebox{0.9}{\tikzfig{realcliffordaxioms/Z2CNOTZ2}} = \scalebox{0.9}{\tikzfig{realcliffordaxioms/CNOTZ1}}
            \end{equation}
        \end{minipage}\vspace{-0.6em}
        \begin{minipage}[c]{0.32\linewidth}\centering
            \begin{equation}\tag{CZr}\label{realXC}
                \scalebox{0.9}{\tikzfig{realcliffordaxioms/HHCNOTHH}} = \scalebox{0.9}{\tikzfig{realcliffordaxioms/NOTC}}
            \end{equation}
        \end{minipage}
        %%%%%% 3-wire %%%%%%
        \begin{minipage}[c]{0.57\linewidth}\centering
            \begin{equation}\tag{CF}\label{realCF}
                \scalebox{0.9}{\tikzfig{realcliffordaxioms/CZcommitCX-left}} = \scalebox{0.9}{\tikzfig{realcliffordaxioms/CZcommitCX-right}}
            \end{equation}
        \end{minipage}
        \begin{minipage}[c]{0.29\linewidth}\centering
            \begin{equation}\tag{I}\label{realInew}
                \scalebox{0.9}{\tikzfig{qutritcliffordaxioms/CNOT23CNOT12CNOT13}} = \scalebox{0.9}{\tikzfig{qutritcliffordaxioms/CNOT12CNOT23}}
            \end{equation}
        \end{minipage}
    \end{minipage}}
    \caption{Real Clifford presentation \(\QCrealclifford\), where the real phase, Hadamard-\(Z\), controlled-\(X\), controlled-\(Z\), and three-wire interaction laws replace the corresponding Clifford rules in the real fragment.}
    \label{fig:qubitrealcliffordaxioms}
\end{figure}
%%%%%%%%%%%%%%%%%%%%%%%%%%%%%%%%%%%%%%%%%%%%%%%%%%%%%%%%%%%%%%%%%%%%%%%%%%%%%%%%%%%%%
%%%%%%%%%%%%%%%%%%%%%%%%%%%%%%%%%%% QUTRIT CLIFFORD %%%%%%%%%%%%%%%%%%%%%%%%%%%%%%%%%%
%%%%%%%%%%%%%%%%%%%%%%%%%%%%%%%%%%%%%%%%%%%%%%%%%%%%%%%%%%%%%%%%%%%%%%%%%%%%%%%%%%%%%
\begin{figure}[ht]
    \centering
    \fbox{\begin{minipage}{0.975\linewidth}
    \centering
    \vspace{-1em}
    \begin{minipage}[c]{0.24\linewidth}\centering
      \begin{equation}\tag{\omegarulabel{12}}\label{qt-omegapow12}
        \scalebox{0.9}{\tikzfig{gates/wb}}^{\otimes 12} = \scalebox{0.9}{\tikzfig{gates/empty}}
      \end{equation}
    \end{minipage}
    \begin{minipage}[c]{0.29\linewidth}\centering
      \begin{equation}\tag{\rulabel{H^4}}\label{qt-hpow4}
        \scalebox{0.9}{\tikzfig{qutritcliffordaxioms/hadpow4}} = \scalebox{0.9}{\tikzfig{cliffordaxioms/I}}
      \end{equation}
    \end{minipage}
    \begin{minipage}[c]{0.27\linewidth}\centering
      \begin{equation}\tag{\rulabel{S^3}}\label{qt-spow3}
        \scalebox{0.9}{\tikzfig{qutritcliffordaxioms/spow3}} = \scalebox{0.9}{\tikzfig{gates/Id}}
      \end{equation}
    \end{minipage}\vspace{-0.6em}
    \begin{minipage}[c]{0.58\linewidth}\centering
      \begin{equation}\tag{E}\label{qt-shpow3}
        \scalebox{0.9}{\tikzfig{qutritcliffordaxioms/shs}} = \scalebox{0.9}{\tikzfig{qutritcliffordaxioms/h3s2h3}}
      \end{equation}
    \end{minipage}\vspace{-0.6em}
    \begin{minipage}[c]{0.43\linewidth}\centering
      \begin{equation}\tag{\rulabel{SS'}}\label{qt-ssprime}
        \scalebox{0.9}{\tikzfig{qutritcliffordaxioms/ssprime}} = \scalebox{0.9}{\tikzfig{qutritcliffordaxioms/sprimes}}
      \end{equation}
    \end{minipage}
    \begin{minipage}[c]{0.35\linewidth}\centering
      \begin{equation}\tag{CPh}\label{qt-cnotremove}
        \scalebox{0.9}{\tikzfig{qutritcliffordaxioms/CNOT-S-K-CNOT}} = \scalebox{0.9}{\tikzfig{qutritcliffordaxioms/S-K}}
      \end{equation}
    \end{minipage}\vspace{-0.6em}
    \begin{minipage}[c]{0.38\linewidth}\centering
      \begin{equation}\tag{KC}\label{qt-kcnot}
        \scalebox{0.9}{\tikzfig{qutritcliffordaxioms/K-CNOT}} = \scalebox{0.9}{\tikzfig{qutritcliffordaxioms/CNOT-CNOT-K}}
      \end{equation}
    \end{minipage}
    \begin{minipage}[c]{0.57\linewidth}\centering
      \begin{equation}\tag{CZ}\label{qt-czdecomp}
        \scalebox{0.9}{\tikzfig{qutritcliffordaxioms/CZ}} = \scalebox{0.9}{\tikzfig{qutritcliffordaxioms/CZdecomp}}
      \end{equation}
    \end{minipage}\vspace{-0.6em}
    \begin{minipage}[c]{0.36\linewidth}\centering
      \begin{equation}\tag{B}\label{qt-swapdecomp}
        \scalebox{0.9}{\tikzfig{cliffordaxioms/SWAPCNOT}} = \scalebox{0.9}{\tikzfig{qutritcliffordaxioms/CNOT-NOTC-NOTC-K}}
      \end{equation}
    \end{minipage}
    % 3-wire
    \begin{minipage}[c]{0.29\linewidth}\centering
      \begin{equation}\tag{I}\label{qt-I}
        \scalebox{0.9}{\tikzfig{qutritcliffordaxioms/CNOT23CNOT12CNOT13}} = \scalebox{0.9}{\tikzfig{qutritcliffordaxioms/CNOT12CNOT23}}
      \end{equation}
    \end{minipage}
    \end{minipage}}
    \caption{Qutrit Clifford presentation \(\QCqutritclifford\), using the qutrit scalar convention and the qutrit versions of the phase, controlled-addition, controlled-\(Z\), swap-decomposition, and three-wire laws.}
    \label{fig:qutritcliffordaxioms}
\end{figure}
%%%%%%%%%%%%%%%%%%%%%%%%%%%%%%%%%%%%%%%%%%%%%%%%%%%%%%%%%%%%%%%%%%%%%%%%%%%%%%%%%%%%%%%
%%%%%%%%%%%%%%%%%%%%%%%%%%%%%%%%%%% QUBIT CLIFFORD+T %%%%%%%%%%%%%%%%%%%%%%%%%%%%%%%%%%
%%%%%%%%%%%%%%%%%%%%%%%%%%%%%%%%%%%%%%%%%%%%%%%%%%%%%%%%%%%%%%%%%%%%%%%%%%%%%%%%%%%%%%%
\begin{figure}[ht]
    \centering
    \fbox{\begin{minipage}{0.975\textwidth}
        \centering
        \vspace{-1em}
        %%%%%% 0-wire %%%%%%%
        \begin{minipage}[c]{0.22\textwidth}
            \centering
            \begin{equation}\tag{\omegarulabel{8}}\label{t-w8}
                \scalebox{0.9}{\tikzfig{gates/wc}}^{\otimes 8} = \scalebox{0.9}{\tikzfig{gates/empty}}
            \end{equation}
        \end{minipage}
        %%%%%% 1-wire %%%%%%%
        \begin{minipage}[c]{0.46\textwidth}
            \centering
            \begin{equation}\tag{\rulabel{T^8}}\label{t-Tpow8}
                \scalebox{0.9}{\tikzfig{cliffordplustaxioms/Tpow8}} = \scalebox{0.9}{\tikzfig{gates/Id}}
            \end{equation}
        \end{minipage}\vspace{-0.6em}
        \begin{minipage}[c]{0.24\textwidth}
            \centering
            \begin{equation}\tag{\rulabel{H^2}}\label{t-H2}
                \scalebox{0.9}{\tikzfig{cliffordaxioms/HH}} = \scalebox{0.9}{\tikzfig{gates/Id}}
            \end{equation}
        \end{minipage}
        \begin{minipage}[c]{0.52\textwidth}
            \centering
            \begin{equation}\tag{E}\label{t-E}
                \scalebox{0.9}{\tikzfig{cliffordaxioms/HSH}} = \scalebox{0.9}{\tikzfig{cliffordaxioms/wSdaggerHSdagger}}
            \end{equation}
        \end{minipage}\vspace{-0.6em}
        \begin{minipage}[c]{0.34\textwidth}
            \centering
            \begin{equation}\tag{TX}\label{t-TX}
                \scalebox{0.9}{\tikzfig{cnotdihedral/R11L}}
            =
            \scalebox{0.9}{\tikzfig{cnotdihedral/R11R}}
          \end{equation}
        \end{minipage}
        %%%%%% 2-wire %%%%%%%
        \begin{minipage}[c]{0.34\textwidth}
            \centering
            \begin{equation}\tag{CPh}\label{t-Ct}
                \scalebox{0.9}{\tikzfig{cliffordplustaxioms/CNOTTCNOT}} = \scalebox{0.9}{\tikzfig{cliffordplustaxioms/T}}
            \end{equation}
        \end{minipage}\vspace{-0.6em}
        \begin{minipage}[c]{0.29\textwidth}
            \centering
            \begin{equation}\tag{B}\label{t-B}
                \scalebox{0.9}{\tikzfig{cliffordaxioms/CNOTNOTC}} = \scalebox{0.9}{\tikzfig{cliffordaxioms/SWAPCNOT}}
            \end{equation}
        \end{minipage}
        \begin{minipage}[c]{0.43\textwidth}
            \centering
            \begin{equation}\tag{CZ}\label{t-CZ}
                \scalebox{0.9}{\tikzfig{cliffordaxioms/HCNOTH}} = \scalebox{0.9}{\tikzfig{cliffordaxioms/CZ}}
            \end{equation}
        \end{minipage}\vspace{-0.6em}
        \begin{minipage}[c]{0.32\textwidth}
            \centering
            \begin{equation}\tag{CSH}\label{t-CsCh}
                \scalebox{0.9}{\tikzfig{cliffordplustaxioms/CsCh}} = \scalebox{0.9}{\tikzfig{cliffordplustaxioms/ChCs}}
            \end{equation}
        \end{minipage}
        \begin{minipage}[c]{0.50\textwidth}
            \centering
            \begin{equation}\tag{\rulabel{HT^2}}\label{t-HTHT}
                \scalebox{0.9}{\tikzfig{cliffordplustaxioms/HTHT}} = \scalebox{0.9}{\tikzfig{cliffordplustaxioms/THTH}}
            \end{equation}
        \end{minipage}\vspace{-0.6em}
        \begin{minipage}[c]{0.65\textwidth}
            \centering
            \begin{equation}\tag{HTH}\label{t-C20}
                \scalebox{0.9}{\tikzfig{cliffordplustaxioms/C20l}} = \scalebox{0.9}{\tikzfig{cliffordplustaxioms/C20r}}
            \end{equation}
        \end{minipage}
    \end{minipage}}
    \caption{Qubit Clifford\(+T\) presentation \(\QCcliffordplust\), combining the Clifford core with the \(T\)-phase, controlled-phase, controlled-Hadamard, and \(HTH\)-type relations needed for the two-qubit completeness transfer.}
    \label{fig:qubitcliffordplustaxioms}
\end{figure}
%%%%%%%%%%%%%%%%%%%%%%%%%%%%%%%%%%%%%%%%%%%%%%%%%%%%%%%%%%%%%%%%%%%%%%%%%%%%%%%%%%%%%%%%
%%%%%%%%%%%%%%%%%%%%%%%%%%%%%%%%%%% QUBIT CLIFFORD+CS %%%%%%%%%%%%%%%%%%%%%%%%%%%%%%%%%%
%%%%%%%%%%%%%%%%%%%%%%%%%%%%%%%%%%%%%%%%%%%%%%%%%%%%%%%%%%%%%%%%%%%%%%%%%%%%%%%%%%%%%%%%
\begin{figure}[ht]
    \centering
    \fbox{\begin{minipage}{0.975\textwidth}
        \centering
        \vspace{-1em}
        %%%%%% 0-wire %%%%%%%
        \begin{minipage}[c]{0.22\textwidth}
            \centering
            \begin{equation}\tag{\omegarulabel{8}}\label{cs-w8}
                \scalebox{0.9}{\tikzfig{gates/wc}}^{\otimes 8} = \scalebox{0.9}{\tikzfig{gates/empty}}
            \end{equation}
        \end{minipage}
        %%%%%% 1-wire %%%%%%%
        \begin{minipage}[c]{0.24\textwidth}
            \centering
            \begin{equation}\tag{\rulabel{H^2}}\label{cs-H2}
                \scalebox{0.9}{\tikzfig{cliffordaxioms/HH}} = \scalebox{0.9}{\tikzfig{gates/Id}}
            \end{equation}
        \end{minipage}
        \begin{minipage}[c]{0.32\textwidth}
            \centering
            \begin{equation}\tag{\rulabel{S^4}}\label{cs-S4}
                \scalebox{0.9}{\tikzfig{cliffordaxioms/SSSS}} = \scalebox{0.9}{\tikzfig{gates/Id}}
            \end{equation}
        \end{minipage}\vspace{-0.6em}
        \begin{minipage}[c]{0.51\textwidth}
            \centering
            \begin{equation}\tag{E}\label{cs-E}
                \scalebox{0.9}{\tikzfig{cliffordaxioms/HSH}} = \scalebox{0.9}{\tikzfig{cliffordaxioms/wSdaggerHSdagger}}
            \end{equation}
        \end{minipage}\vspace{-0.6em}
        %%%%%% 2-wire %%%%%%%
        \begin{minipage}[c]{0.41\textwidth}
            \centering
            \begin{equation}\tag{CPh}\label{cs-C}
                \scalebox{0.9}{\tikzfig{cliffordpluscsaxioms/CSpow4withS}} = \scalebox{0.9}{\tikzfig{cliffordaxioms/S}}
            \end{equation}
        \end{minipage}
        \begin{minipage}[c]{0.31\textwidth}
            \centering
            \begin{equation}\tag{CSr}\label{cs-CSrev}
                \scalebox{0.9}{\tikzfig{cliffordpluscsaxioms/CS}} = \scalebox{0.9}{\tikzfig{cliffordpluscsaxioms/CSrev}}
            \end{equation}
        \end{minipage}\vspace{-0.6em}
        \begin{minipage}[c]{0.28\textwidth}
            \centering
            \begin{equation}\tag{B}\label{cs-B}
                \scalebox{0.9}{\tikzfig{cliffordaxioms/CNOTNOTC}} = \scalebox{0.9}{\tikzfig{cliffordaxioms/SWAPCNOT}}
            \end{equation}
        \end{minipage}
        \begin{minipage}[c]{0.59\textwidth}
            \centering
            \begin{equation}\tag{XCS}\label{cs-XCS}
                \scalebox{0.9}{\tikzfig{cliffordpluscsaxioms/X2TopLeft}}=\scalebox{0.9}{\tikzfig{cliffordpluscsaxioms/X2TopRight}}
            \end{equation}
        \end{minipage}\vspace{-0.6em}
        \begin{minipage}[c]{0.49\textwidth}
            \centering
            \begin{equation}\tag{CE}\label{cs-SHCHC}
                \scalebox{0.9}{\tikzfig{cliffordpluscsaxioms/HTopLeft}}=\scalebox{0.9}{\tikzfig{cliffordpluscsaxioms/HTopRight}}
            \end{equation}
        \end{minipage}
        %%%%%%% 3-wire %%%%%%%
        \begin{minipage}[c]{0.29\textwidth}
            \centering
            \begin{equation}\label{cs-I}\tag{I}\scalebox{0.9}{\tikzfig{qutritcliffordaxioms/CNOT23CNOT12CNOT13}} = \scalebox{0.9}{\tikzfig{qutritcliffordaxioms/CNOT12CNOT23}}\end{equation}
        \end{minipage}\vspace{-0.6em}
        \begin{minipage}[c]{0.64\textwidth}
            \centering
            \begin{equation}\label{cs-U}\tag{\rulabel{SH_k}}\scalebox{0.9}{\tikzfig{cliffordpluscsaxioms/CSmonsterL}} = \scalebox{0.9}{\tikzfig{cliffordpluscsaxioms/CSmonsterR}}\end{equation}
        \end{minipage}
    \end{minipage}}
    \caption{Qubit Clifford\(+CS\) presentation \(\QCcliffordpluscs\), combining the Clifford phase laws with controlled-\(S\) commutation and the two three-wire relations used for the three-qubit transfer; in \cref{cs-U}, the dashed box repeats the enclosed subcircuit \(k\) times for \(k\in\{0,1,2,3\}\).}
    \label{fig:qubitcliffordpluscsaxioms}
\end{figure}
%%%%%%%%%%%%%%%%%%%%%%%%%%%%%%%%%%%%%%%%%%%%%%%%%%%%%%%%%%%%%%%%%%%%%%%%%%%%%%%%%%%%%%%%%%
%%%%%%%%%%%%%%%%%%%%%%%%%%%%%%%%%%% QUBIT CNOT-DIHEDRAL %%%%%%%%%%%%%%%%%%%%%%%%%%%%%%%%%%
%%%%%%%%%%%%%%%%%%%%%%%%%%%%%%%%%%%%%%%%%%%%%%%%%%%%%%%%%%%%%%%%%%%%%%%%%%%%%%%%%%%%%%%%%%
\begin{figure}[ht]
    \centering
    \fbox{%
        \begin{minipage}{0.98\textwidth}
        \centering
        \vspace{-1em}
        \begin{minipage}[c]{0.22\textwidth}
            \centering
            \begin{equation}\tag{\omegarulabel{8}}\label{cnot-R10}
                \scalebox{0.9}{\tikzfig{gates/wc}}^{\otimes 8} = \scalebox{0.9}{\tikzfig{gates/empty}}
            \end{equation}
        \end{minipage}
        \begin{minipage}[c]{0.27\textwidth}
            \centering
            \begin{equation}\tag{\rulabel{X^2}}\label{cnot-R1}
            \scalebox{0.9}{\tikzfig{cnotdihedral/R1L}}
            =
            \scalebox{0.9}{\tikzfig{gates/Id}}
            \end{equation}
        \end{minipage}
        \begin{minipage}[c]{0.24\textwidth}
            \centering
            \begin{equation}\tag{\rulabel{T^8}}\label{cnot-R7}
            \scalebox{0.9}{\tikzfig{cnotdihedral/R7L}}
            =
            \scalebox{0.9}{\tikzfig{gates/Id}}
            \end{equation}
        \end{minipage}\vspace{-0.6em}
        \begin{minipage}[c]{0.34\textwidth}
            \centering
            \begin{equation}\tag{TX}\label{cnot-R11}
            \scalebox{0.9}{\tikzfig{cnotdihedral/R11L}}
            =
            \scalebox{0.9}{\tikzfig{cnotdihedral/R11R}}
            \end{equation}
        \end{minipage}
        \begin{minipage}[c]{0.34\textwidth}
            \centering
            \begin{equation}\tag{XC}\label{new-R3}
            \scalebox{0.9}{\tikzfig{identities/wCNOT-00}}
            =
            \scalebox{0.9}{\tikzfig{identities/wCNOT-16}}
            \end{equation}
        \end{minipage}
        \begin{minipage}[c]{0.29\textwidth}
            \centering
            \begin{equation}\tag{B}\label{new-R5}
            \scalebox{0.9}{\tikzfig{cliffordaxioms/CNOTNOTC}}
            =
            \scalebox{0.9}{\tikzfig{cliffordaxioms/SWAPCNOT}}
            \end{equation}
        \end{minipage}\vspace{-0.6em}
        \begin{minipage}[c]{0.35\textwidth}
            \centering
            \begin{equation}\tag{ZC}\label{new-R8}
            \scalebox{0.9}{\tikzfig{identities/zCNOT-00}}
            =
            \scalebox{0.9}{\tikzfig{identities/zCNOT-05}}
            \end{equation}
        \end{minipage}
        \begin{minipage}[c]{0.29\textwidth}
            \centering
            \begin{equation}\tag{CPh}\label{cnot-R12}
            \scalebox{0.9}{\tikzfig{cliffordplustaxioms/CNOTTCNOT}}
            =
            \scalebox{0.9}{\tikzfig{cliffordplustaxioms/T}}
            \end{equation}
        \end{minipage}\vspace{-0.6em}
        \begin{minipage}[c]{0.29\textwidth}
            \centering
            \begin{equation}\tag{I}\label{cnot-R6}
            \scalebox{0.9}{\tikzfig{cnotdihedral/R6L}}
            =
            \scalebox{0.9}{\tikzfig{cnotdihedral/R6R}}
            \end{equation}
        \end{minipage}
        \begin{minipage}[c]{0.48\textwidth}
            \centering
            \begin{equation}\tag{\rulabel{C^2T}}\label{cnot-R9}
            \scalebox{0.9}{\tikzfig{cnotdihedral/R9L}}
            =
            \scalebox{0.9}{\tikzfig{cnotdihedral/R9R}}
            \end{equation}
        \end{minipage}\vspace{-0.6em}
        \begin{minipage}[c]{0.84\textwidth}
            \centering
            \begin{equation}\tag{\rulabel{C^3T}}\label{cnot-R13}
            \scalebox{0.82}{\tikzfig{cnotdihedral/R13L}}
            =
            \scalebox{0.82}{\tikzfig{cnotdihedral/R13R}}
            \end{equation}
        \end{minipage}
        \end{minipage}%
    }
    \caption{CNOT-dihedral presentation \(\QCdihedral\), with the \(X\), \(T\), CNOT, phase-commutation, \(C^2T\), and \(C^3T\) relations used to recover the omitted source axioms.}
    \label{fig:qubitdihedralaxioms}
\end{figure}
